\scanpage{34}
Courbes elliptiques sur un \deleted{anneau de} corps local à car. résiduelle $p>0$: étude infinitésimale.

$A$ anneau de val. discrète complet, corps des fractions $K$, corps résiduel $k$ de car $p>0$.

$E_K$ courbe elliptique sur $K$, \deleted{\unsure{d'où un grp Module|d'où un grp Modèle}}
\figure{\begin{tikzpicture}\draw[->] (0,0.3) -- (0,0);\end{tikzpicture}}

1º) \uline{$K$ de car 0} i.e. $T_p(E_{\overline{K}})_{\text{ét}}$ \uline{de rang 2} \deleted{\unsure{libre}}.
\margin{\unsure{[}et $T_p(E_{\overline{K}})_{\text{ét}}=T_p(E_{\overline{K}})$]}

On veut étudier les structures possibles pour ce module galoisien.

a) $E_K$ a ``bonne réduction'' i.e. $E$ provient d'un schéma elliptique $E$ sur $A$, \uline{et} $E_k$ est ``ordinaire''.

Alors $T_p(E_K)$ est \inserted{can.} extension d'un \deleted{module} \inserted{$L$} de rang 1, \inserted{\uline{m. n. ramifié}}, par un \deleted{autre} module de la forme
\[
\deleted{\operatorname{Hom}}\quad L^{\vee}\otimes T_p(G_m)=L^{\vee}[1].
\]
\margin{\diagram{\begin{tikzcd}[ampersand replacement=\&,column sep=large]
0\arrow[r]\&L^{\vee}(1)\arrow[r,phantom,"{\unsure{\text{---}|\longrightarrow}}" description]\&T_p(E_K)\arrow[r]\&L^{\unsure{*|\vee}}\arrow[r,phantom,"{\unsure{\text{---}|\longrightarrow}}" description]\&\unsure{\cdot|0}
\end{tikzcd}}}
De plus, l'extension, \deleted{qui} est donnée par un élément de
\[
H^1(K,\deleted{\unsure{L^{\vee}\otimes L^{\vee}|\operatorname{Hom}(L,L^{\vee})}}L^{-2}\otimes_{Z_p}T_p(G_m)),
\]
qui dans le cas où $A$ est à corps résiduel sép. clos (donc $L$ trivial sur $A$) s'identifie à
\[
H^1(K,T_p(G_m))\otimes_{Z_p}L_0^{-2}
\simeq\left(\varprojlim\unsure{K^*/\deleted{\unsure{Z|\otimes}}{\inserted{K^*}}^{p^n}|K^*\otimes Z/p^n}\right)\otimes_{Z_p}L_0^{-2}
\simeq\operatorname{Hom}(L_0^2,\widehat{K^*}).
\]
\margin{$L_0$: valeur en $k$}

\margin{33}
